% 子文件：E_涂色分数
% 本文件包含 6 个 section



%% ==============================
%% 小数的意义：通过图形涂色理解小数
%% ==============================
\newpage

\section{解答：小数的意义与图形涂色}

\vspace{0.4cm}

\noindent\textbf{下面每个图形都表示整数“1”。选择合适的图形，分别涂色表示 0.009、0.35、0.4。}

\vspace{0.8cm}

\begin{center}
\begin{tikzpicture}[>=Stealth, scale=1.3] % 整体思路：三幅图分别表示十分之几、百分之几和千分之几；先定义统一网格样式，再分别画竖条图、百格图和立体千格图，并把对应部分涂黑

  % 统一定义网格线样式
  \tikzset{ % 在当前图内临时定义网格线与边框样式，便于三幅图统一风格
    gridline/.style={cyan!80!black, line width=0.4pt}, % gridline 表示内部细网格线：青黑色、线宽 0.4pt
    frame/.style={cyan!80!black, line width=0.8pt} % frame 表示外框线：同色但更粗，线宽 0.8pt
  } % 结束局部样式定义

  % =========== 第1个图：10个竖条 (0.4) ===========
  \begin{scope}[shift={(0,0)}] % 第一幅图放在左侧
    % 涂色 (0.4 对应 4 个竖条)
    \fill[black!50] (0,0) rectangle (0.8, 2); % 把左边 4 个竖条涂黑；0.8=4×0.2，对应 0.4
    
    % 竖线
    \foreach \i in {1,...,9} { % 画 9 条内部竖分割线，把整体分成 10 个等宽竖条
      \draw[gridline] (\i*0.2, 0) -- (\i*0.2, 2); % 第 i 条竖线位于 x=0.2i 处，从 y=0 画到 y=2
    } % 结束第一幅图竖线循环
    % 外框
    \draw[frame] (0,0) rectangle (2, 2); % 画第一幅图外框，整体表示整数 1
    
    \node[below=15pt, font=\large] at (1, 0) {( \textcolor{blue}{0.4} )}; % 在图下方写对应小数 0.4
  \end{scope} % 结束第一幅“十分之几”图

  % =========== 第2个图：10x10格子 (0.35) ===========
  \begin{scope}[shift={(3.8,0)}] % 第二幅图右移 3.8cm
    % 涂色 (0.35 对应 35 个小格, 也就是上方30格加第4排的5格)
    \fill[black!50] (0, 1.4) rectangle (2, 2); % 先涂满上面 3 行，共 30 个小格
    \fill[black!50] (0, 1.2) rectangle (1.0, 1.4); % 再涂第 4 行左边 5 格，共计 35 个小格，对应 0.35
    
    % 内部网格
    \draw[gridline, step=0.2cm] (0,0) grid (2, 2); % 用步长 0.2cm 画 10×10 百格网
    % 外框
    \draw[frame] (0,0) rectangle (2, 2); % 画百格图外框
    
    \node[below=15pt, font=\large] at (1, 0) {( \textcolor{blue}{0.35} )}; % 在图下方写 0.35
  \end{scope} % 结束第二幅“百分之几”图

  % =========== 第3个图：10x10x10正方体 (0.009) ===========
  \begin{scope}[shift={(7.8,-0.2)}] % 第三幅立体图右移并略下移
    \def\dx{0.8} % 定义立体图向右后方的水平偏移量 dx
    \def\dy{0.45} % 定义立体图向右后方的竖直偏移量 dy
    
    % 涂色 (0.009 对应前面板底侧的 9 个小格)
    \fill[black!50] (0, 0) rectangle (1.8, 0.2); % 在正面底部涂 9 个小方格；1.8=9×0.2，表示 0.009
    
    % 前面板内部网格
    \foreach \i in {1,...,9} { % 在正面画横竖网格，把正方形面分成 10×10 小格
      \draw[gridline] (0, \i*0.2) -- (2, \i*0.2); % 第 i 条横线
      \draw[gridline] (\i*0.2, 0) -- (\i*0.2, 2); % 第 i 条竖线
    } % 结束正面网格循环
    \draw[frame] (0,0) rectangle (2,2); % 画正面外框
    
    % 上面板内部网格
    \foreach \i in {1,...,9} { % 在上面板画透视网格
      \draw[gridline] (\i*0.08, 2+\i*0.045) -- (2+\i*0.08, 2+\i*0.045); % 画与前边平行的上表面横向透视线
      \draw[gridline] (\i*0.2, 2) -- (\i*0.2+\dx, 2+\dy); % 画从前上边缘延伸到后上边缘的斜向分割线
    } % 结束上面板网格循环
    \draw[frame] (0,2) -- (\dx, 2+\dy) -- (2+\dx, 2+\dy) -- (2,2) -- cycle; % 画上面板外框
    
    % 右边板内部网格
    \foreach \i in {1,...,9} { % 在右侧面继续画透视网格
      \draw[gridline] (2, \i*0.2) -- (2+\dx, \i*0.2+\dy); % 画右侧面的斜向横格线
      \draw[gridline] (2+\i*0.08, \i*0.045) -- (2+\i*0.08, 2+\i*0.045); % 画右侧面的竖向透视线
    } % 结束右侧面网格循环
    \draw[frame] (2,0) -- (2+\dx, \dy) -- (2+\dx, 2+\dy) -- (2,2) -- cycle; % 画右侧面外框
    
    \node[below=15pt, font=\large] at (1+\dx/2, 0) {( \textcolor{blue}{0.009} )}; % 在第三幅图下方写 0.009
  \end{scope} % 结束第三幅“千分之几”立体图

\end{tikzpicture}
\end{center}

\vspace{0.6cm}
{\small\color{gray}
\textbf{解析：}\\
1. \textbf{0.4}：一位小数表示十分之几。左侧第一个图形被平均分成了 $10$ 个竖条，给其中的 $4$ 个竖条涂上颜色，即表示整体的十分之四（也就是 $0.4$）。\\
2. \textbf{0.35}：两位小数表示百分之几。中间的图形被分成了 $10 \times 10 = 100$ 个小方格，给其中的 $35$ 个方格涂上颜色，即表示整体的百分之三十五（也就是 $0.35$）。通常从上面开始，整齐地涂满 $3$ 行（$30$ 格），再加上第 $4$ 行的前 $5$ 格。\\
3. \textbf{0.009}：三位小数表示千分之几。右侧的立体图形被分作了 $10 \times 10 \times 10 = 1000$ 个小正方体。取其中的 $9$ 个涂上颜色，即占整个大正方体的千分之九（也就是 $0.009$）。只需挑选其正面底排最左侧连续的 $9$ 个小方格涂色即可。
}

\vspace{0.6cm}

{\small\color{blue!70!black}
\textbf{绘图基本原理与步骤：}\\
\textbf{1. 多图并排布局}：使用 \texttt{scope} 环境搭配 \texttt{shift} 平移算子，将三幅独立图形（十分位竖条图、百分位百格图、千分位立体图）横向排列于同一 \texttt{tikzpicture} 内，统一管理缩放系数 \texttt{scale=1.3}。\\
\textbf{2. 涂色区域精确计算}：涂色面积通过 \texttt{fill[black!50]} 的 \texttt{rectangle} 坐标严格按份数换算。如 0.4 对应涂满左起 $4 \times 0.2 = 0.8$ 宽度；0.35 拆分为上方满 3 行加第 4 行前 5 格。\\
\textbf{3. 立体透视模拟}：第三幅千格图通过定义偏移量 \texttt{dx=0.8, dy=0.45} 构建伪三维正方体的上面板和右侧面，利用平行四边形框架和透视网格线模拟立体效果。
}






%% ==============================
%% 分数的初步认识：涂色表示分数
%% ==============================
\newpage

\section{解答：在图中涂色表示分数}

\vspace{0.4cm}

\noindent\textbf{1. 在图中涂色表示它下面的分数。}

\vspace{0.8cm}

\begin{center}
\begin{tikzpicture}[>=Stealth, scale=0.85, line width=1pt] % 整体思路：三幅图分别通过“平均分后涂色”的方式表示分数，因此每一幅都先画外形或网格，再填充对应份数，最后在下方标分数

  % 1. 正方形网格 3x3
  \begin{scope}[shift={(0,0)}] % 第一幅图保持在左侧原位
    \fill[black!70] (0, 2) rectangle (3, 3); % 先把最上面一整行 3 个小正方形涂黑；rectangle 表示用对角点画矩形
    \fill[black!70] (0, 1) rectangle (2, 2); % 再把中间一行左边 2 个小格涂黑，这样一共涂了 5 个小格
    \draw[line width=0.8pt, step=1cm] (0,0) grid (3,3); % 画 3×3 网格；step=1cm 表示每一格边长 1cm；line width=0.8pt 为细边框
    \node[below=15pt, font=\Large] at (1.5, 0) {$\frac{5}{9}$}; % 在图下方 15pt 处写出分数 5/9
  \end{scope} % 结束第一幅正方形网格图

  % 2. 平行四边形
  \begin{scope}[shift={(5.5,0)}] % 第二幅图整体右移 5.5cm
    % 涂色区域 (左上角子平行四边形)
    \fill[black!70] (0.75, 1.5) -- (1.95, 1.5) -- (2.7, 3) -- (1.5, 3) -- cycle; % 用四个顶点围成一个小平行四边形并填黑，表示 4 份中的 1 份；cycle 自动闭合图形
    % 外框
    \draw[line width=0.8pt] (0,0) -- (2.4,0) -- (3.9,3) -- (1.5,3) -- cycle; % 画大平行四边形外框
    % 分割线
    \draw[line width=0.8pt] (1.2,0) -- (2.7,3); % 画一条与侧边平行的分割线，把图形纵向分成两份
    \draw[line width=0.8pt] (0.75,1.5) -- (3.15,1.5); % 再画一条水平样式的分割线，把图形横向分成两份，最终得到 4 等份
    % 分数
    \node[below=15pt, font=\Large] at (1.95, 0) {$\frac{1}{4}$}; % 在图下方标出对应分数 1/4
  \end{scope} % 结束第二幅平行四边形分数图

  % 3. 椭圆与圆圈
  \begin{scope}[shift={(13.5, 1.5)}] % 第三幅图右移并略上移，让椭圆与前两图对齐
    % 椭圆外框
    \draw[line width=0.8pt] (0, 0) ellipse (3.6cm and 1.6cm); % 以原点为中心画椭圆；横半轴 3.6cm，纵半轴 1.6cm，表示整体“1”的容器
    
    % 排列圆圈 (2行 x 5列)
    \foreach \x in {-2, -1, 0, 1, 2} { % 外层循环控制 5 列圆圈的位置
      \foreach \y in {-0.6, 0.6} { % 内层循环控制上下两行圆圈的位置
        \ifnum \x<1 % 条件判断：当前列编号若小于 1，就属于左边 3 列，需要涂黑
          % 涂黑的圆圈 (左边 3 列共 6 个)
          \filldraw[fill=black!85, draw=black, line width=0.8pt] (\x*1.1, \y) circle (0.35cm); % filldraw 同时画轮廓并填充；圆心坐标由 \x*1.1 和 \y 给出；circle (0.35cm) 表示半径 0.35cm
        \else
          % 空白圆圈 (右边 2 列共 4 个)
          \draw[line width=0.8pt] (\x*1.1, \y) circle (0.35cm); % 对右边两列只画空心圆，不填色
        \fi
      }
    }
    % 分数
    \node[below=15pt, font=\Large] at (0, -1.5) {$\frac{3}{5}$}; % 在椭圆下方写出分数 3/5
  \end{scope} % 结束第三幅“圆圈分组”分数图

\end{tikzpicture}
\end{center}

\vspace{0.6cm}
{\small\color{gray}
\textbf{解析：}\\
1. $\frac{5}{9}$：把正方形平均分成 $9$ 份，涂色部分占其中的 $5$ 份（涂满上方的 $3$ 个和中间行左侧的 $2$ 个小正方形）。\\
2. $\frac{1}{4}$：把平行四边形平均分成 $4$ 份，涂色部分占其中的 $1$ 份（本题涂满左上方的 $1$ 个平行四边形）。\\
3. $\frac{3}{5}$：共有 $10$ 个圆圈，平均分作 $5$ 份，每份 $2$ 个。$\frac{3}{5}$ 表示取其中的 $3$ 份，即 $2 \times 3 = 6$ 个圆圈，故将左侧的 $6$ 个小圆圈全部涂成黑色实体。
}

\vspace{0.8cm}


\vspace{0.6cm}

{\small\color{blue!70!black}
\textbf{绘图基本原理与步骤：}\\
\textbf{1. 坐标定位与几何构建}：根据题意中的已知长度和角度，使用 \texttt{coordinate} 定义关键顶点坐标，通过 \texttt{draw} 指令按序连接成完整几何图形。线宽、颜色参数区分原图（黑色）与答案（红色 \texttt{red!70!black}）图层。\\
\textbf{2. 辅助量标注}：利用 \texttt{node} 的方向锚点（\texttt{above}、\texttt{below}、\texttt{left}、\texttt{right}）将角度、边长、面积等数据标注在图形周围，避免与线条或其他标注重叠。若涉及角弧则使用 \texttt{arc} 或 \texttt{pic[angle]} 命令渲染。\\
\textbf{3. 虚实线分层}：题干已知条件用实线绘制，解答新增的辅助线或操作痕迹用 \texttt{dashed} 虚线或彩色加粗线标识，确保读者一眼区分原有信息与作答内容。
}

%% ==============================
%% 分数、小数、百分数互化
%% ==============================
\newpage




\newpage

\section{解答：十分格与百分格的小数涂色}

\vspace{0.4cm}

\noindent\textbf{3. 每个大正方形都表示整数“$1$”，涂色表示下面的小数。}

\vspace{0.8cm}

\begin{center}
\begin{tikzpicture}[>=Stealth, x=0.5cm, y=0.5cm]

  % =========== 左侧：0.3 (十分之三) ===========
  % 构建一个巨大的物理尺寸为 5cm x 5cm 的主正方形块
  \begin{scope}[shift={(0,0)}]
    % 1. 先铺设底层色彩
    % 左侧图像为被平分为 10 个横条，表示 0.3 需要涂 3 等份。
    % 这里模拟原图逻辑，将顶部 3 行 (y=7至y=10 跨度) 高亮为蓝色。
    \fill[blue!55!cyan, opacity=0.85] (0, 7) rectangle (10, 10);
    
    % 2. 绘制中层内部网格线 (防止被加重颜色的底色遮盖)
    \foreach \y in {1,2,...,9} {
      \draw[thin, black!70] (0,\y) -- (10,\y);
    }
    
    % 3. 绘制强制覆盖的最外周黑框极黑粗线
    \draw[thick, black] (0,0) rectangle (10,10);
    
    % 4. 底部题注说明标签
    \node[below] at (5, -0.6) {\Large $0.3$};
  \end{scope}

  % =========== 右侧：0.42 (百分之四十二) ===========
  % 中心向右偏移 14 格（7cm），以确保充裕的呼吸缓冲留白。
  \begin{scope}[shift={(14,0)}]
    % 1. 组合式区块底层铺色
    % 因原图要求精确显示 42 个百格区块，拒绝笨重的循环，直接分两组大区块矢量渲染：
    % - 前 4 整列矩形，宽 4 格，高满 10 (共 40 面积单元)
    \fill[blue!55!cyan, opacity=0.85] (0, 0) rectangle (4, 10);
    % - 位于第 5 列的地基零散碎片，宽 1 格，仅包含底部 2 层高度 (即 2 面积单元)
    \fill[blue!55!cyan, opacity=0.85] (4, 0) rectangle (5, 2);
    
    % 2. 覆盖并悬浮于颜色色块之上的精密虚线方网格骨架
    % 这种分次渲染的方式真实营造了“彩色方块格子”的物理立体纵深感！
    \draw[step=1, thin, black!70] (0,0) grid (10,10);
    
    % 3. 外沿加粗强切边保护
    \draw[thick, black] (0,0) rectangle (10,10);
    
    % 4. 底部题注说明标签
    \node[below] at (5, -0.6) {\Large $0.42$};
  \end{scope}

\end{tikzpicture}
\end{center}

\vspace{0.8cm}

{\small\color{gray}
\textbf{解析：}\\
1. \textbf{小数 $0.3$ 的理解与涂色}：左侧的大正方形被单向地分作了 $10$ 根平行的相等长横条。在此模型中每一横条所占整个正方形的比例即表示为整体“$1$”的 $\frac{1}{10}$，即 $0.1$。要直观展示 $0.3$，题境即需要我们在图中均匀抹色 $3$ 根这样的横条。为了几何呈现的美感与系统合规的连贯性，在参考解答中，程序逻辑选择将排于最先导位置的顶端 $3$ 条（从全局视点俯视涵盖了矩阵中高度向最顶层的 $30\%$）统一执行蓝色面曝光涂层。\\
2. \textbf{小数 $0.42$ 的理解与涂色}：右侧的大正方形矩阵环境显得截然不同，它由纵横均匀切割的蛛网线均分为了 $10 \times 10 = 100$ 个极微小正方形单位格。至此，宏观环境下的每一个零次小方格均代表着总体“$1$”的 $\frac{1}{100}$，即 $0.01$；而如果纵向贯穿集齐其中整个一条列柱（包扩内部无缝衔接的 $10$ 个底层方格微元区）也就聚合成了规模为 $0.1$ 的高阶表示。为了完美组装出 $0.42$ 这个数值，我们可以运用位运算思想将其拆分归类为 $0.4 + 0.02$。因此，画笔自左侧基点发端，优先向右扩张抹满前 $4$ 个大型纵队整列（其面积精确等价于 $0.4$，涵盖了多达 $40$ 枚小微元矩阵格），紧接着，其笔锋滑坠至第 $5$ 列底部地基点突击向上仅继续漫灌充盈 $2$ 枚独立孤格（这一组残段等同于 $0.02$ 的占格），所有图形单元总和面积精密啮合形成绝妙的 $42$ 格几何图斑！
}

\vspace{0.6cm}

{\small\color{blue!70!black}
\textbf{绘图基本原理与步骤：}\\
\textbf{1. 几何原语的无损解耦分层技术}：底层绘制程序高度奉行教科书插画级的严格 Z-Index 时序切分制（Z-排序深度叠加技术）。将充满着自带高抗锯齿防混叠通道配方的 \texttt{blue!55!cyan} 深潜色块优先排版至图形序列中的原生态级底层 \texttt{fill} 域进行纯色烘培；尔后在这个高强色彩膜壳表面启动独立线程，进行后续那些代表着切分的网格与切线痕迹的高像素扫描刻画工艺 \texttt{draw}；保证任何细若游丝的横向骨架虚实缝隙都不因色彩漫溢被意外剥离吞噬，进而将题目要求的“只涂满方格内的区域空间感”的拟真要求刻勒得一览无遗。\\
\textbf{2. 百分格的极简时空多边形拼接工艺}：针对右侧要求极高分辨率下呈现出惊人的 $42\%$ 细碎网点涂色矩阵区，原生计算核心否绝了那些依靠冗散代码逐个填充微格从而让程序极度负载并可能崩溃的嵌套大循环算式路线；相反地，程序在幕后构建了一场堪称绝妙的多边形大型矩形无缝熔焊拼接战术。仅仅凭靠一道指令下发了长宽量级极其极端的纵向大跨度坐标区间矢量矩形 \texttt{[0,4] $\times$ [0,10]}，用以瞬间一笔拉满左侧重器极的 $40$ 格基石壁垒；其后追加一枚微型补充探针指令，打在了第 $5$ 列地表浅区 \texttt{[4,5] $\times$ [0,2]}。巨集模块与微格碎片之间的边缘参数精确自锁咬合，未生半点撕裂缝合毛刺！
}


\newpage



%% ==============================
%% 第91题（补充题）：六边形的分数涂色
%% ==============================
\newpage

\section{解答：六边形的分数涂色}

\vspace{0.6cm}

\textbf{\LARGE 82.} \textbf{将每个六边形分别涂色表示下面的分数。\\[1ex]
第一图下标注：$\displaystyle\frac{1}{2}$ \qquad 第二图下标注：$\displaystyle\frac{1}{3}$ \qquad 第三图下标注：$\displaystyle\frac{1}{6}$}

\vspace{0.4cm}
\textbf{解：}

\textbf{分析：}
如图所示，每个正六边形都被内部相交的三条对角线平均分成了 $6$ 个完全相同的小正三角形。
要表示特定的分数，只需要计算应涂色的小三角形数量：
1. 第一图表示 $\displaystyle\frac{1}{2}$：需要涂色总块数的 $\displaystyle\frac{1}{2}$，即 $6 \times \frac{1}{2} = 3$ 块。
2. 第二图表示 $\displaystyle\frac{1}{3}$：需要涂色总块数的 $\displaystyle\frac{1}{3}$，即 $6 \times \frac{1}{3} = 2$ 块。
3. 第三图表示 $\displaystyle\frac{1}{6}$：需要涂色总块数的 $\displaystyle\frac{1}{6}$，即 $6 \times \frac{1}{6} = 1$ 块。

\vspace{0.6cm}
\begin{center}
\begin{tikzpicture}[>=Stealth, scale=1.6]

  % 定义绘制一个被6等分的正六边形的宏
  % 参数1: x位移
  \def\drawHexagon#1{
    \begin{scope}[shift={(#1, 0)}]
      % 画六边形外框
      \draw[line width=1pt] (0:1) \foreach \i in {1,2,3,4,5} { -- (\i*60:1) } -- cycle;
      % 画三条虚线对角线
      \draw[dashed, line width=0.6pt] (0:1) -- (180:1);
      \draw[dashed, line width=0.6pt] (60:1) -- (240:1);
      \draw[dashed, line width=0.6pt] (120:1) -- (300:1);
    \end{scope}
  }

  % 第一个图形：涂 3 块 (1/2)
  \begin{scope}[shift={(0,0)}]
    % 填充 0~180 度的 3 个小正三角形
    \fill[blue!30] (0,0) -- (0:1) \foreach \i in {1,2,3} { -- (\i*60:1) } -- cycle;
    \drawHexagon{0}
    \node[below, font=\Large] at (0, -1.2) {$\displaystyle\frac{1}{2}$};
  \end{scope}

  % 第二个图形：涂 2 块 (1/3)
  \begin{scope}[shift={(3.5,0)}]
    % 填充 60~180 度的 2 个小正三角形 (顶部及左上)
    \fill[blue!30] (0,0) -- (60:1) \foreach \i in {2,3} { -- (\i*60:1) } -- cycle;
    \drawHexagon{0}
    \node[below, font=\Large] at (0, -1.2) {$\displaystyle\frac{1}{3}$};
  \end{scope}

  % 第三个图形：涂 1 块 (1/6)
  \begin{scope}[shift={(7,0)}]
    % 填充 60~120 度的 1 个小正三角形 (正上方)
    \fill[blue!30] (0,0) -- (60:1) -- (120:1) -- cycle;
    \drawHexagon{0}
    \node[below, font=\Large] at (0, -1.2) {$\displaystyle\frac{1}{6}$};
  \end{scope}

\end{tikzpicture}
\end{center}

\vspace{0.4cm}
{\small\color{blue!70!black}
\textbf{绘图原理与步骤：}\\
\textbf{1. 极坐标正多边形生成：} 利用极坐标系循环生成法 \texttt{(A:R)} 快速渲染完全标准正六边形，设定外接圆定长半径为 $1$。并人为规定顶点的极角依次为 $0^\circ, 60^\circ, 120^\circ \dots 300^\circ$。这一设定不仅确保了图形的顶边与底边绝对水平，且使得连接 $0^\circ$ 和 $180^\circ$ 的对角线刚好是一根完美的水平中轴线，与原题素材的物理放置姿态完全锲合。\\
\textbf{2. 均等基本块划分：} 通过追加 \texttt{\textbackslash draw[dashed]} 指令连接全部核心对角线，实现在渲染外框的同时，在图形内部打通了隔断线。物理上将原本致密的整块大正六边形成功均分为 $6$ 块等面积的正三角形阵列单元。\\
\textbf{3. 分割填色法则：} 将题目提出的代数分数逻辑映射为涂色数量逻辑（$1/2\to 3$ 块，$1/3\to 2$ 块，$1/6\to 1$ 块）。然后启用特定的区块填充层 \texttt{\textbackslash fill[blue!30]}，精确约束中心 $(0,0)$ 延伸到外围对应节点之间的扇形发散式多边形面域，彻底杜绝手工填充导致的色彩溢出或留白死角。
}

%% ==============================
%% 第92题（补充题）：用不同方法表示正方形的 1/4
%% ==============================
\newpage



%% ==============================
%% 第92题（补充题）：用不同方法表示正方形的 1/4
%% ==============================
\newpage

\section{解答：表示分数的多种折叠模型}

\vspace{0.6cm}

\textbf{\LARGE 83.} \textbf{（第 4 题）你能用不同的方法表示每个正方形的 $\displaystyle\frac{1}{4}$ 吗？先用纸折一折，再画一画。}

\vspace{0.4cm}
\textbf{解：}

\textbf{分析：}
要把正方形平均分成 $4$ 份并着色表示出其中的 $\displaystyle\frac{1}{4}$，常见的基于对称性的折法（等分法）共有以下 4 种典型图案：
1. \textbf{竖向等分}：将正方形往同一方向连续对折两次，得到 4 个相同高度的竖直长条。涂色其中 1 块。
2. \textbf{横向等分}：与竖向等分同理，得到 4 个相同宽度的横向长条。涂色其中 1 块。
3. \textbf{十字等分}：横向和纵向各中心对折一次，得到 4 个相同的小正方形。涂色其中 1 块小正方形。
4. \textbf{对角线等分}：沿着正方形的两条对角线分别对折，得到 4 个相同的等腰直角三角形。涂色其中 1 块。

\vspace{0.6cm}
\begin{center}
\begin{tikzpicture}[>=Stealth, x=2.2cm, y=2.2cm]

  % 提取绘制单体正方形的宏
  \def\drawSquare#1{
    % 画实线外框
    \draw[line width=1pt] (#1, 0) rectangle (#1+1, 1);
  }

  % 1. 竖向四等分
  \begin{scope}[shift={(0,0)}]
    % 涂色第一条
    \fill[blue!20] (0, 0) rectangle (0.25, 1);
    \drawSquare{0}
    % 画分划虚线
    \draw[dashed, line width=0.6pt] (0.25, 0) -- (0.25, 1);
    \draw[dashed, line width=0.6pt] (0.5,  0) -- (0.5,  1);
    \draw[dashed, line width=0.6pt] (0.75, 0) -- (0.75, 1);
    \node[below, font=\small] at (0.5, -0.25) {方法一};
  \end{scope}

  % 2. 纵向四等分（横向长条）
  \begin{scope}[shift={(1.4,0)}]
    % 涂色最下一条
    \fill[blue!20] (0, 0) rectangle (1, 0.25);
    \drawSquare{0}
    \draw[dashed, line width=0.6pt] (0, 0.25) -- (1, 0.25);
    \draw[dashed, line width=0.6pt] (0, 0.5)  -- (1, 0.5);
    \draw[dashed, line width=0.6pt] (0, 0.75) -- (1, 0.75);
    \node[below, font=\small] at (0.5, -0.25) {方法二};
  \end{scope}

  % 3. 十字四等分
  \begin{scope}[shift={(2.8,0)}]
    % 涂色左下角小正方形
    \fill[blue!20] (0, 0) rectangle (0.5, 0.5);
    \drawSquare{0}
    \draw[dashed, line width=0.6pt] (0.5, 0) -- (0.5, 1);
    \draw[dashed, line width=0.6pt] (0, 0.5) -- (1, 0.5);
    \node[below, font=\small] at (0.5, -0.25) {方法三};
  \end{scope}

  % 4. 对角线四等分
  \begin{scope}[shift={(4.2,0)}]
    % 涂色下方三角形
    \fill[blue!20] (0, 0) -- (1, 0) -- (0.5, 0.5) -- cycle;
    \drawSquare{0}
    \draw[dashed, line width=0.6pt] (0, 0) -- (1, 1);
    \draw[dashed, line width=0.6pt] (0, 1) -- (1, 0);
    \node[below, font=\small] at (0.5, -0.25) {方法四};
  \end{scope}

\end{tikzpicture}
\end{center}

\vspace{0.4cm}
{\small\color{blue!70!black}
\textbf{绘图原理与步骤：}\\
\textbf{1. 图元定位阵列：} 建立坐标域，定义基础正方形内部比例常数为 $1$。分别于 $x=0, 1.4, 2.8, 4.2$ 的偏移处并行铺设四个基准大实体正方形外框，还原原题干中横向并排依次排开的四大方框环境。\\
\textbf{2. 几何代数等分模型：} 展现物理操作的“折纸”效果即是向框架内追加平分级的辅助线。针对等分这一严密的数学命题，分别编写了三等距绝对垂直律、绝对水平分割律、正交十字对偶中心律以及点对称极限折痕（交叉对角）等四种拓扑方程，通过 \texttt{dashed} 属性赋予拆分折痕。\\
\textbf{3. 独立集块高亮着色：} 利用 \texttt{\textbackslash fill[blue!20]} 的特定区域填充功能，对于矩形切块严格采用 \texttt{rectangle} 参数矩阵锁死，对于三角切块则采用闭环多边形连线 \texttt{cycle} 限位，强行高亮出四种形态截然不同但面积恒等占比 $1/4$ 的比例单元。
}


%% ==============================
%% 第93题（补充题）：分数涂色与大小比较
%% ==============================
\vspace{1.5cm}

\textbf{\LARGE 84.} \textbf{（第 3 题）先涂色表示分数，再比较每组分数的大小。}

\vspace{0.4cm}
\textbf{解：}

\textbf{分析：}
比较给定的两组真分数：
1. 第一组：同分母分数 $\displaystyle\frac{3}{8}$ 与 $\displaystyle\frac{5}{8}$ 比较。分母相同代表切分的标准块大小一样，此时“分子越大，分数实际所占的面积就越大”，故 $\displaystyle\frac{3}{8} < \displaystyle\frac{5}{8}$。
对应在图上的视觉直观操作：每个大正方形已被细分为 $8$ 个等分小三角形，在左图中涂去 $3$ 个，在右图中涂去 $5$ 个，清晰可见右侧涂色区更大。
2. 第二组：同分母分数 $\displaystyle\frac{3}{4}$ 与 $\displaystyle\frac{2}{4}$ 比较。同理，由分子数值可判明 $\displaystyle\frac{3}{4} > \displaystyle\frac{2}{4}$。
在图上的动作：每个大正方形已被切分为 $4$ 个等分小方格，左图按分子占涂 $3$ 个，右图涂 $2$ 个即见大小高低。

\vspace{0.6cm}
\begin{center}
\begin{tikzpicture}[>=Stealth, x=1.8cm, y=1.8cm]

  %% 第一组：3/8 与 5/8
  \begin{scope}[shift={(0,0)}]
    % 左边正方形 (3/8)
    \begin{scope}[shift={(0,0)}]
      % 涂 3 个小三角形
      \fill[blue!30] (0.5, 0.5) -- (0,0) -- (0.5,0) -- cycle;
      \fill[blue!30] (0.5, 0.5) -- (0.5,0) -- (1,0) -- cycle;
      \fill[blue!30] (0.5, 0.5) -- (1,0) -- (1,0.5) -- cycle;
      
      % 外框
      \draw[line width=1pt] (0,0) rectangle (1,1);
      % 米字分割线
      \draw[dashed, line width=0.6pt] (0.5,0) -- (0.5,1);
      \draw[dashed, line width=0.6pt] (0,0.5) -- (1,0.5);
      \draw[dashed, line width=0.6pt] (0,0) -- (1,1);
      \draw[dashed, line width=0.6pt] (0,1) -- (1,0);
    \end{scope}

    % 右边正方形 (5/8)
    \begin{scope}[shift={(1.4,0)}]
      % 涂 5 个小三角形
      \fill[blue!30] (0.5, 0.5) -- (0,0) -- (0.5,0) -- cycle;
      \fill[blue!30] (0.5, 0.5) -- (0.5,0) -- (1,0) -- cycle;
      \fill[blue!30] (0.5, 0.5) -- (1,0) -- (1,0.5) -- cycle;
      \fill[blue!30] (0.5, 0.5) -- (1,0.5) -- (1,1) -- cycle;
      \fill[blue!30] (0.5, 0.5) -- (1,1) -- (0.5,1) -- cycle;
      
      % 外框
      \draw[line width=1pt] (0,0) rectangle (1,1);
      % 米字分割线
      \draw[dashed, line width=0.6pt] (0.5,0) -- (0.5,1);
      \draw[dashed, line width=0.6pt] (0,0.5) -- (1,0.5);
      \draw[dashed, line width=0.6pt] (0,0) -- (1,1);
      \draw[dashed, line width=0.6pt] (0,1) -- (1,0);
    \end{scope}

    % 下方分数及比较圆圈符号
    \node[anchor=mid, font=\Large] at (0.5, -0.6) {$\displaystyle\frac{3}{8}$};
    \draw[line width=0.8pt] (1.2, -0.6) circle (0.2);
    \node[anchor=mid, font=\large] at (1.2, -0.6) {$<$};
    \node[anchor=mid, font=\Large] at (1.9, -0.6) {$\displaystyle\frac{5}{8}$};
  \end{scope}


  %% 第二组：3/4 与 2/4
  \begin{scope}[shift={(4.5,0)}]
    % 左边正方形 (3/4)
    \begin{scope}[shift={(0,0)}]
      % 涂 3 个小正方形
      \fill[blue!30] (0,0) rectangle (1,0.5);
      \fill[blue!30] (0,0.5) rectangle (0.5,1);
      
      % 外框
      \draw[line width=1pt] (0,0) rectangle (1,1);
      % 十字分割线
      \draw[dashed, line width=0.6pt] (0.5,0) -- (0.5,1);
      \draw[dashed, line width=0.6pt] (0,0.5) -- (1,0.5);
    \end{scope}

    % 右边正方形 (2/4)
    \begin{scope}[shift={(1.4,0)}]
      % 涂 2 个小正方形
      \fill[blue!30] (0,0) rectangle (1,0.5);
      
      % 外框
      \draw[line width=1pt] (0,0) rectangle (1,1);
      % 十字分割线
      \draw[dashed, line width=0.6pt] (0.5,0) -- (0.5,1);
      \draw[dashed, line width=0.6pt] (0,0.5) -- (1,0.5);
    \end{scope}

    % 下方分数及比较圆圈符号
    \node[anchor=mid, font=\Large] at (0.5, -0.6) {$\displaystyle\frac{3}{4}$};
    \draw[line width=0.8pt] (1.2, -0.6) circle (0.2);
    \node[anchor=mid, font=\large] at (1.2, -0.6) {$>$};
    \node[anchor=mid, font=\Large] at (1.9, -0.6) {$\displaystyle\frac{2}{4}$};
  \end{scope}

\end{tikzpicture}
\end{center}

\vspace{0.4cm}
{\small\color{blue!70!black}
\textbf{绘图原理与步骤：}\\
\textbf{1. 域隔离对比系：} 利用两个独立包裹的 \texttt{scope} （平移量相隔足够宽距）宏观分区建立了两套绝对独立的对照绘图区。左侧为八等分体系，右侧为基准十字四等分体系，保持所有方块外延尺寸严密一致并排布局，彻底实现原题左右互不干扰的并排比列版面。\\
\textbf{2. 细分网格裁切场：} 对于“涂色表示分数的实质计算”，必须对空白基板做等效切分：左侧区域通过对角线加横竖中位线簇（施加 \texttt{dashed} 点划线属性）切割成 $8$ 个“全边等底正三角”；右侧区域则仅走单纯的十字正中切分分离为 $4$ 个纯正方微块块，为底层的数据结构绑定图元。\\
\textbf{3. 分区面积映射法则：} 紧扣两组代数算子中的分子系数（$3、5$ 及 $3、2$），指令程序采用天蓝色（\texttt{blue!30}）精确检索并吞噬掉对应个数的微型切块。并在最后一行布置底注级渲染——用中心坐标 \texttt{anchor=mid} 精准压平实线评判空圆圈，同时往圆心中注入机器级准确判断的 $<$ 及 $>$ 等比例排版的矢量不等号。
}

%% ==============================
%% 第94题（补充题）：等值分数涂色
%% ==============================
\newpage



%% ==============================
%% 第94题（补充题）：等值分数涂色
%% ==============================
\newpage

\section{解答：等值分数的涂色与转化}

\vspace{0.6cm}

\textbf{\LARGE 85.} \textbf{先在图中涂色表示已知分数的等值分数，再填空。}

\vspace{0.4cm}
\textbf{解：}

\textbf{分析：}
“等值分数”是指虽然分子分母的数字不同，但其表示的实质数值（在一个固定大小的图形中即指代对应相等的覆盖面积）永远相等的分数。
1. 第一个图形组：左侧的六边形被极点射线分成了 $3$ 个并排的菱形，阴影部分占据 $1$ 个菱形，即表示 $\displaystyle\frac{1}{3}$。为了保证面积的绝对“等值”，右侧的六边形（已被全等切分为 $6$ 个小正三角形）其涂色区域必须与左图一致（占据左半边），此时刚好吞并了 $2$ 块小三角形，故得出等价分数公式 $\displaystyle\frac{1}{3} = \frac{2}{6}$。
2. 第二个图形组：左侧的正方形横向等分切成了 $4$ 份水平层，顶层被涂成了灰色，客观表示 $\displaystyle\frac{1}{4}$。而在右侧的正方形中，整个画面被进一步碎化成 $4 \times 4 = 16$ 个同尺寸的小方格，要想使得涂色面积与左边高度吻合一致，那就必须涂满第一排平行的 $4$ 个小方格，由此得出 $\displaystyle\frac{1}{4} = \frac{4}{16}$。

\vspace{0.6cm}
\begin{center}
\begin{tikzpicture}[>=Stealth, line width=0.8pt]

  %% ========== 1. 左侧六边形组 (1/3 和 2/6) ==========
  \begin{scope}[shift={(0,0)}]
    \def\R{1.4}
    % 左一：涂灰色 (题干要求)
    \fill[gray!30] (0,0) -- (120:\R) -- (180:\R) -- (240:\R) -- cycle;
    % 边界底侧六边形
    \draw (0:\R) \foreach \i in {1,...,5} { -- (\i*60:\R) } -- cycle;
    % 内部三连平分虚线
    \draw[dashed, line width=0.6pt] (0,0) -- (0:\R);
    \draw[dashed, line width=0.6pt] (0,0) -- (120:\R);
    \draw[dashed, line width=0.6pt] (0,0) -- (240:\R);
  \end{scope}

  % 左侧等式推导： 1/3 = 2/6
  \node[anchor=mid, font=\Large] at (1.8, -2.1) {$\displaystyle\frac{1}{3} \quad = \quad \frac{(\quad \textcolor{blue!80!black}{2} \quad)}{(\quad \textcolor{blue!80!black}{6} \quad)}$};

  \begin{scope}[shift={(3.6,0)}]
    \def\R{1.4}
    % 左二：涂蓝色作为解答图
    \fill[blue!30] (0,0) -- (120:\R) -- (180:\R) -- (240:\R) -- cycle;
    \draw (0:\R) \foreach \i in {1,...,5} { -- (\i*60:\R) } -- cycle;
    % 对角交叉虚线 (六等分)
    \draw[dashed, line width=0.6pt] (0:\R) -- (180:\R);
    \draw[dashed, line width=0.6pt] (60:\R) -- (240:\R);
    \draw[dashed, line width=0.6pt] (120:\R) -- (300:\R);
  \end{scope}


  %% ========== 2. 右侧正方形组 (1/4 和 4/16) ==========
  \begin{scope}[shift={(7.5,-1.2)}]
    \def\S{2.4}
    % 右一：涂灰色 (最顶层1/4)
    \fill[gray!30] (0, 0.75*\S) rectangle (\S, \S);
    \draw (0,0) rectangle (\S, \S);
    \draw[dashed, line width=0.6pt] (0, 0.25*\S) -- (\S, 0.25*\S);
    \draw[dashed, line width=0.6pt] (0, 0.5*\S)  -- (\S, 0.5*\S);
    \draw[dashed, line width=0.6pt] (0, 0.75*\S) -- (\S, 0.75*\S);
  \end{scope}

  % 右侧等式推导： 1/4 = 4/16
  \node[anchor=mid, font=\Large] at (10.6, -2.1) {$\displaystyle\frac{1}{4} \quad = \quad \frac{(\quad \textcolor{blue!80!black}{4} \quad)}{(\quad \textcolor{blue!80!black}{16} \quad)}$};

  \begin{scope}[shift={(11.3,-1.2)}]
    \def\S{2.4}
    % 右二：涂天蓝色解答 (最上面一排共计 4 个小格子)
    \fill[blue!30] (0, 0.75*\S) rectangle (\S, \S);
    \draw (0,0) rectangle (\S, \S);
    % 内部密织 16 宫格虚线阵列
    \foreach \i in {0.25, 0.5, 0.75} {
      \draw[dashed, line width=0.6pt] (0, \i*\S) -- (\S, \i*\S);
      \draw[dashed, line width=0.6pt] (\i*\S, 0) -- (\i*\S, \S);
    }
  \end{scope}

\end{tikzpicture}
\end{center}

\vspace{0.4cm}
{\small\color{blue!70!black}
\textbf{绘图原理与步骤：}\\
\textbf{1. 极坐标正方多边形衍生：} 沿袭运用完美的极坐标生成宏构建底侧的基础正六边形，结合虚线平分技术实现 $1/3$ 到 $1/6$ 的精密切割过渡；在此环节中，原题提供基准图被统一施配灰色 \texttt{gray!30} 透明图层充当环境前提背景。\\
\textbf{2. 欧几里得网格切分场：} 根据正方形切片映射原理，我们对相同基底尺寸（宽幅 $2.4\text{ cm}$）的正方形采取了两套割裂动作逻辑。一是单纯运用 \texttt{y} 轴刻度线实施粗颗粒切划；二是使用了二维复套循环 \texttt{\textbackslash foreach} 将其进行极限十六级粉碎化制备。\\
\textbf{3. 核心解题反馈印记：} 在待作答的图形轮廓上严格用“主观解答色——淡蓝色 \texttt{blue!30}”追溯描摹出了完全等比的等积投影版块（譬如那完美的 $4$ 枚连排小方格）；与此同时借由底部的等式渲染结构，极其自然地将 $2/6$ 及 $4/16$ 利用代码着色一并填入圆润的小括号之中。
}


